% ════════════════════════════════════════════════════════════════
%  主要符号与记号约定（frontmatter，全局统一符号）
%  注：本章在 frontmatter（chapter 计数器=0），故不用 \caption（会编成“表 0.x”），
%      各组用居中粗体小标题 + booktabs 表格呈现。各章末“符号表”是本表的局部投影。
% ════════════════════════════════════════════════════════════════
\chapter*{主要符号与记号约定}
\addcontentsline{toc}{chapter}{主要符号与记号约定}
\label{ch:notation}
\markboth{主要符号与记号约定}{主要符号与记号约定}

本章汇总全书统一的记号与核心约定，供随时查阅；各章末的“符号表”是本表在该章的局部投影。

\section*{核心约定}
\begin{itemize}
  \item \textbf{字体}：标量用普通体 $a,\lambda$；向量用粗体小写 $\mathbf{a},\boldsymbol{\phi}$（默认列向量）；矩阵用粗体大写 $\mathbf{A},\mathbf{R}$；集合与流形用花体 $\mathcal{M},\mathcal{X}$；李代数用哥特体 $\so(3),\se(3)$。
  \item \textbf{帧与位姿记号}：$\mathbf{R}_{ab}$、$\mathbf{t}_{ab}$ 把 $b$ 系下的坐标变换到 $a$ 系，即 $\mathbf{p}_a=\mathbf{R}_{ab}\,\mathbf{p}_b+\mathbf{t}_{ab}$（下标读作“$a\leftarrow b$”）；例如 $\mathbf{R}_{wc}$ 把相机系下的坐标变换到世界系。\emph{注意}：各文献对上下标方向的约定不一——例如 $\mathbf{R}_a^b$ 在有的书中表示 $a\to b$（即本书的 $\mathbf{R}_{ba}$，与本书相反），而 ${}^a\mathbf{R}_b$ 在有的书中表示 $b\to a$（与本书相同）；援引他书结论前务必先核对其定义。
  \item \textbf{扰动}：流形（$\SO(3)/\SE(3)$）上的求导与更新\textbf{以右扰动为主线}（局部坐标，便于对接因子图与 GTSAM/Ceres）；涉及左扰动的文献结论以并列对照给出。
  \item \textbf{四元数}：采用 \textbf{Hamilton} 约定，实部在前 $\mathbf{q}=[\,w,\ \mathbf{v}\,]$；$\mathbf{q}$ 与 $-\mathbf{q}$ 表示同一旋转（双重覆盖）。
  \item \textbf{$\se(3)$ 坐标}：取“平移在前” $\boldsymbol{\xi}=[\boldsymbol{\rho};\boldsymbol{\phi}]\in\mathbb{R}^6$（$\boldsymbol{\rho}$ 平移部分、$\boldsymbol{\phi}$ 旋转部分；一般 $\boldsymbol{\rho}\neq\mathbf{t}$，见 $\SE(3)$ 指数映射）。
\end{itemize}

\section*{符号一览}

\begin{center}\small
\textbf{通用与线性代数}\par\vspace{3pt}
\begin{tabular}{ll}
\toprule
记号 & 含义 \\
\midrule
$a,\ \lambda$ & 标量 \\
$\mathbf{a},\ \boldsymbol{\phi}$ & 向量（粗体小写，默认列向量） \\
$\mathbf{A},\ \mathbf{R},\ \mathbf{T}$ & 矩阵（粗体大写） \\
$\mathbf{I}_n$ & $n\times n$ 单位矩阵（维数明确时略下标） \\
$(\cdot)^\top,\ (\cdot)^{-1}$ & 转置 / 逆 \\
$\lVert\cdot\rVert,\ \langle\cdot,\cdot\rangle$ & 范数 / 内积 \\
$\mathbf{a}\times\mathbf{b}=\mathbf{a}^\wedge\mathbf{b}$ & 三维叉积（$(\cdot)^\wedge$ 见下表） \\
$\det(\cdot),\ \operatorname{tr}(\cdot)$ & 行列式 / 迹 \\
\bottomrule
\end{tabular}
\end{center}

\begin{center}\small
\textbf{旋转、位姿、李群与李代数}\par\vspace{3pt}
\begin{tabular}{ll}
\toprule
记号 & 含义 \\
\midrule
$\SO(3),\ \SE(3)$ & 三维旋转群 / 刚体位姿群 \\
$\so(3),\ \se(3)$ & 对应的李代数 \\
$\mathbf{R}\in\SO(3),\ \mathbf{T}\in\SE(3)$ & 旋转矩阵 / 齐次变换矩阵 \\
$\mathbf{R}_{ab},\ \mathbf{t}_{ab}$ & $b$ 系到 $a$ 系的旋转 / 平移 \\
$(\cdot)^\wedge,\ (\cdot)^\vee$ & 向量 $\leftrightarrow$ 李代数（$\so(3)$ 反对称化；$\se(3)$ 重载） \\
$\exp,\ \log$ & 矩阵指数 / 对数 \\
$\mathrm{Exp},\ \mathrm{Log}$ & 吃/吐向量的指数 / 对数（$\mathrm{Exp}(\boldsymbol{\phi})=\exp(\boldsymbol{\phi}^\wedge)$） \\
$\boldsymbol{\phi}=\theta\mathbf{a}$ & 旋转向量（轴角：$\theta$ 转角、$\mathbf{a}$ 单位轴） \\
$\boldsymbol{\xi}=[\boldsymbol{\rho};\boldsymbol{\phi}]$ & $\se(3)$ 坐标（平移在前） \\
$\mathbf{J}_l,\ \mathbf{J}_r$ & 左 / 右雅可比（$3\times3$） \\
$\boldsymbol{\mathcal{J}}_l$ & $\SE(3)$ 的 $6\times6$ 左雅可比 \\
$\mathrm{Ad}(\mathbf{T}),\ \mathrm{ad}(\boldsymbol{\xi})$ & 群伴随（$6\times6$）/ 小伴随 \\
$(\cdot)^\odot$ & 齐次点扰动算子（$4\times6$，$\SE(3)$ 右扰动） \\
$\boxplus,\ \boxminus$ & 流形上的“加 / 减”（retraction / 局部坐标） \\
$\mathbf{q}=[\,w,\mathbf{v}\,]$ & 单位四元数（Hamilton，实部在前） \\
\bottomrule
\end{tabular}
\end{center}

\begin{center}\small
\textbf{概率、估计与优化}\par\vspace{3pt}
\begin{tabular}{ll}
\toprule
记号 & 含义 \\
\midrule
$\mathbb{E}[\cdot],\ \operatorname{Cov}[\cdot]$ & 期望 / 协方差 \\
$\mathcal{N}(\boldsymbol{\mu},\boldsymbol{\Sigma})$ & 高斯分布（均值 $\boldsymbol{\mu}$、协方差 $\boldsymbol{\Sigma}$） \\
$\boldsymbol{\Sigma}_{\mathbf{w}},\ \boldsymbol{\Sigma}_{\mathbf{v}}$ & 过程噪声 / 观测噪声协方差 \\
$\boldsymbol{\Lambda},\ \boldsymbol{\eta}$ & 信息矩阵（$=\boldsymbol{\Sigma}^{-1}$；部分章节记作 $\boldsymbol{\Omega}$，等价）/ 信息向量 \\
$\check{\mathbf{x}},\ \hat{\mathbf{x}}$ & 先验估计 / 后验估计 \\
$\lVert\mathbf{e}\rVert_{\boldsymbol{\Sigma}}^2=\mathbf{e}^\top\boldsymbol{\Sigma}^{-1}\mathbf{e}$ & 马氏（白化）平方范数 \\
$\mathbf{J}=\partial\mathbf{e}/\partial\mathbf{x}$ & 残差雅可比 \\
$\mathbf{H}=\mathbf{J}^\top\mathbf{W}^{-1}\mathbf{J}$ & 高斯–牛顿近似 Hessian（信息矩阵） \\
$\mathbf{W}=\boldsymbol{\Sigma}$ & 权重矩阵（堆叠观测/过程协方差；其逆 $\mathbf{W}^{-1}$ 为信息阵） \\
\bottomrule
\end{tabular}
\end{center}

\begin{center}\small
\textbf{几何与相机}\par\vspace{3pt}
\begin{tabular}{ll}
\toprule
记号 & 含义 \\
\midrule
$\mathbf{K}$ & 相机内参矩阵 \\
$\pi(\cdot),\ \pi^{-1}(\cdot)$ & 投影 / 反投影算子 \\
$\mathbf{P}=[X,Y,Z]^\top$ & 三维空间点 \\
$\mathbf{p}=[u,v]^\top$ & 像素坐标 \\
$b,\ d$ & 双目基线 / 视差（$z=fb/d$） \\
$\tilde{(\cdot)}$ & 齐次坐标 \\
\bottomrule
\end{tabular}
\end{center}
